% !TeX root = ../../main.tex
% sections/app/appendix_aggregate_output.tex

\chapter{集計生産}
\label{chap:main_appendix_aggregate_output}

\section{自国の集計生産}
\label{sec:main_appendix_aggregate_output_home}

\subsection{集計生産\(Y_t^H\)の定義}
\label{subsec:main_appendix_aggregate_output_home_def}
まず生産者物価指数\(p_t^H\)と集計生産\(Y_t^H\)の積が
家計の名目生産の総和\(\sum_{h\in H}p_t^hy_t^h\)となるように集計生産\(Y_t^H\)を定義する。
すなわち
\begin{align}
\label{eq:main_appendix_aggregate_output_home_Y_H_def}
    &p_t^HY_t^H=\sum_{h\in H}p_t^hy_t^h\\
    &\Leftrightarrow Y_t^H=\frac{\sum_{h\in H}p_t^hy_t^h}{p_t^H}
\end{align}

\subsection{家計\(y_t^h\)の財の市場均衡}
\label{sec:main_appendix_aggregate_output_home_equilibrium_y_h}
市場均衡では生産\(y_t^h\)はその財への全世界からの総需要と等しい。
\begin{equation}
\label{eq:main_appendix_aggregate_output_home_equilibrium_y_h_eq}
    y_t^h=\sum_{h'\in H}c_t^{h'\to h}+\sum_{f\in F}c_t^{f\to h}
\end{equation}
ここで右辺の\(c_t^{h'\to h}\)に自国家計\(h'\)の家計\(h\)の財への需要関数
\eqref{eq:model_households_indices_aggregation_stage1a_demand_functions_c_h_h_eq}
を代入し、\(c_t^{f\to h}\)に外国家計\(f\)の家計\(h\)の財への需要関数
\eqref{eq:model_households_indices_aggregation_stage1a_demand_functions_c_f_h_eq}
を代入すると以下のように変形できる。
\begin{align}
\label{eq:main_appendix_aggregate_output_home_equilibrium_y_h_modified_eq}
    y_t^h&=\sum_{h'\in H}\left[\left(\frac{p_t^h}{p_t^H}\right)^{-\theta^H}c_t^{h'\to H}\right]+\sum_{f\in F}\left[\left(\frac{p_t^h}{p_t^H}\right)^{-\theta^H}c_t^{f\to H}\right] \\
    &=\left(\frac{p_t^h}{p_t^H}\right)^{-\theta^H}\left[\sum_{h'\in H}c_t^{h'\to H}+\sum_{f\in F}c_t^{f\to H}\right] \nonumber
\end{align}
ここで大括弧の中身は世界の自国財消費指数の総和を意味する。
これを\(D_t^H\)とおけば以下が得られる。
\begin{equation}
\label{eq:main_appendix_aggregate_output_home_equilibrium_y_h_modified2_eq}
    y_t^h=\left(\frac{p_t^h}{p_t^H}\right)^{-\theta^H}D_t^H
\end{equation}

\subsection{市場均衡における集計生産\(Y_t^H\)と世界の自国財消費指数の総和\(\sum_{h'\in H}c_t^{h'\to H}+\sum_{f\in F}c_t^{f\to H}\)の一致}
\label{sec:main_appendix_aggregate_output_home_Y_H_sum_c_h_H_plus_c_f_H_sum_p_h_y_h}
次に家計の名目生産の総和\(\sum_{h\in H}p_t^hy_t^h\)に上記の財市場均衡における\(y_t^h\)の関係式を代入する。
\begin{align}
\label{eq:main_appendix_aggregate_output_home_Y_H_sum_c_h_H_plus_c_f_H_sum_p_h_y_h_eq}
    \sum_{h\in H}p_t^hy_t^h&=\sum_{h\in H}p_t^h\left[\left(\frac{p_t^h}{p_t^H}\right)^{-\theta^H}D_t^H\right] \\
    &=D_t^H(p_t^H)^{\theta^H}\sum_{h\in H}(p_t^h)^{1-\theta^H}
\end{align}
ここで生産者物価指数\(p_t^H\)の定義\eqref{eq:model_variables_endogenous_aggregate_p_H_def}より
\(\sum_{h\in H}(p_t^h)^{1-\theta^H}=(p_t^H)^{1-\theta^H}\)である。これを代入すると
\begin{align}
\label{eq:main_appendix_aggregate_output_home_Y_H_sum_c_h_H_plus_c_f_H_sum_p_h_y_h_modified_eq}
    \sum_{h\in H}p_t^hy_t^h&=D_t^H(p_t^H)^{\theta^H}(p_t^H)^{1-\theta^H} \\
    &=p_t^HD_t^H
\end{align}
これを\(Y_t^H\)の定義式\(p_t^HY_t^H=\sum p_t^hy_t^h\)と比較すると
集計生産\(Y_t^H\)が世界の自国財消費指数の総和\(D_t^H\)と一致することがわかる。
\begin{equation}
\label{eq:main_appendix_aggregate_output_home_Y_H_sum_c_h_H_plus_c_f_H_Y_H_D_H_eq}
    Y_t^H=D_t^H
\end{equation}

\subsubsection*{市場均衡における集計生産\(Y_t^H\)のCES集計形式による表現}
\label{sec:main_appendix_aggregate_output_home_Y_H_ces}
上記の関係式をもとに集計生産\(Y_t^H\)がCES集計の形式になることを示す。
式\eqref{eq:main_appendix_aggregate_output_home_market_equilibrium_modified2_eq}の両辺を
\(\frac{\theta^H-1}{\theta^H}\)乗し、すべての家計\(h\)について合計をとる。
\begin{align}
\label{eq:main_appendix_aggregate_output_home_Y_H_ces_sum_y_h_eq}
    \sum_{h\in H}(y_t^h)^{\frac{\theta^H-1}{\theta^H}}&=\sum_{h\in H}\left[\left(\frac{p_t^h}{p_t^H}\right)^{-\theta^H}Y_t^H\right]^{\frac{\theta^H-1}{\theta^H}} \\
    &=\frac{(Y_t^H)^{\frac{\theta^H-1}{\theta^H}}}{(p_t^H)^{1-\theta^H}}\sum_{h\in H}(p_t^h)^{1-\theta^H} \nonumber
\end{align}
再び生産者物価指数\(p_t^H\)の定義\eqref{eq:model_variables_endogenous_aggregate_p_H_def}より
\(\sum_{h\in H}(p_t^h)^{1-\theta^H}=(p_t^H)^{1-\theta^H}\)を用いると、右辺の価格項は相殺される。
\begin{equation}
\label{eq:main_appendix_aggregate_output_home_Y_H_ces_sum_y_h_modified_eq}
    \sum_{h\in H}(y_t^h)^{\frac{\theta^H-1}{\theta^H}}=(Y_t^H)^{\frac{\theta^H-1}{\theta^H}}
\end{equation}
この両辺の\(\frac{\theta^H}{\theta^H-1}\)乗を取ると、\(Y_t^H\)が以下のCES集計関数でなければならないことが示される。
\begin{equation}
\label{eq:main_appendix_aggregate_output_home_Y_H_ces_Y_H_eq}
    Y_t^H=\left[\sum_{h\in H}(y_t^h)^{\frac{\theta^H-1}{\theta^H}}\right]^{\frac{\theta^H}{\theta^H-1}}
\end{equation}

\section{外国の集計生産}
\label{sec:main_appendix_aggregate_output_foreign}
同様に外国についても以下が成立する。
外国の集計生産\(Y_t^F\)は生産者物価指数\(p_t^{F*}\)を用いて以下のように定義される。
\begin{align}
\label{eq:main_appendix_aggregate_output_foreign_Y_F_def}
    &p_t^{F^*}Y_t^F=\sum_{f\in F}p_t^{f^*}y_t^f\\
    &\Leftrightarrow Y_t^F=\frac{\sum_{f\in F}p_t^{f^*}y_t^f}{p_t^{F^*}}
\end{align}
市場均衡においてこの集計生産\(Y_t^F\)は世界の外国財消費指数の総和\(D_t^F\coloneqq\sum_{h\in H}c_t^{h\to F}+\sum_{f'\in F}c_t^{f'\to F}\)と一致する。
\begin{equation}
\label{eq:main_appendix_aggregate_output_foreign_Y_F_D_F_eq}
    Y_t^F=D_t^F
\end{equation}
同じく市場均衡において\(Y_t^F\)は以下のCES集計形式により表現される。
\begin{equation}
\label{eq:main_appendix_aggregate_output_foreign_Y_F_ces_Y_F_eq}
    Y_t^F=\left[\sum_{f\in F}(y_t^f)^{\frac{\theta^F-1}{\theta^F}}\right]^{\frac{\theta^F}{\theta^F-1}}
\end{equation}